% Generated draft: requires mathematical and source review.
\documentclass[11pt]{article}
\usepackage[margin=25mm]{geometry}
\usepackage{fontspec}
\setmainfont{Latin Modern Roman}
\usepackage{amsmath,amssymb,mathtools,unicode-math}
\setmathfont{Latin Modern Math}
\usepackage{enumitem,needspace,xurl,hyperref}
\setlength{\parindent}{0pt}
\setlength{\parskip}{5pt}
\setlength{\emergencystretch}{3em}
\newcommand{\N}{\mathbb{N}}
\newcommand{\Z}{\mathbb{Z}}
\newcommand{\Q}{\mathbb{Q}}
\newcommand{\R}{\mathbb{R}}
\newcommand{\C}{\mathbb{C}}
\newcommand{\F}{\mathbb{F}}
\newcommand{\A}{\mathbb{A}}
\newcommand{\PP}{\mathbb{P}}
\newcommand{\E}{\mathbb{E}}
\newcommand{\eps}{\varepsilon}
\newcommand{\dd}{\,\mathrm{d}}
\begin{document}
\section*{UnsolvedMath ID: 10400003}
\textbf{Source problem code:} catalogue entry
\subsection{Definitions and mathematical statement}
For the usual domain of oriented links \(L\subset S^3\), write \(V_L(t)\) for the Jones polynomial. A standard normalization is \(V_{\bigcirc}(t)=1\) and
\[
 t^{-1}V_{L_+}(t)-tV_{L_-}(t)
   =(t^{1/2}-t^{-1/2})V_{L_0}(t),
\]
where the oriented links \(L_+,L_-,L_0\) agree outside a ball and differ inside it by the positive crossing, negative crossing, and oriented smoothing. This convention specifies the notation here; the supplied question itself gives no normalization.

Problem 1.3 asks: find a \(3\)-dimensional topological interpretation of the Jones polynomial of links. This is a request for a construction or characterization, not a stated equality with a specified \(3\)-manifold invariant. The record gives no formal criterion for an acceptable interpretation, no category of links or ambient manifolds beyond the word ``links,'' and no instructions about framings, roots of unity, or recovering the polynomial from evaluations. Those choices must be made explicit by any proposed interpretation rather than inserted as hypotheses into this problem.

The registry labels the item solved and its research summary reports an intended interpretation by \(\mathrm{SU}(2)\) Chern--Simons theory, attributed to Witten (1989), with rigorous associated TQFT construction by Reshetikhin--Turaev. This is the supplied assessment, not a proof or an independently verified statement that every possible meaning of the request is resolved.
\subsection{Short English statement}
Find a topological construction in three dimensions that interprets the Jones polynomial of links. The registry reports a solution in the intended sense through Chern--Simons theory and the Reshetikhin--Turaev TQFT, while the older source presents an open problem. The exact intended sense of interpretation requires source review.
\subsection{Sources}
\begin{itemize}
\item Local UnsolvedMath registry, permanent ID \(10400003\), exact source code \texttt{catalogue entry}, Problem 1.3. The supplied title and statement spell the invariant ``Jon es polynomial''; this draft reads that spacing as Jones polynomial. No individual proposer or proposed year is supplied.
\item Ohtsuki, \emph{Problems on invariants of knots and 3-manifolds} (2002), Problem 1.3, PDF page 10: \url{https://msp.org/gtm/2002/04/gtm-2002-04-024s.pdf}. The monograph's bold headings distinguish 231 numbered problems, conjectures, and questions from surrounding references to them.
\item Supplied extraction: pdf-font-structure; accessed 2026-07-29. Source status NEEDS\_REVIEW says the source presents an open problem and a current primary-literature audit remains. Registry status solved and research classification SOLVED-IN-LITERATURE report Witten (1989), \(\mathrm{SU}(2)\) Chern--Simons theory, and the Reshetikhin--Turaev construction; no full citations for those works are supplied.
\item Dataset attribution: Ulam AI / UnsolvedMath Contributors, CC BY 4.0. Source rights NEEDS\_REVIEW: publicly accessible author/publisher copy, redistribution terms requiring review.
\end{itemize}
\end{document}
